\chapter{Experimentos propuestos}
\section{Medición de la velocidad de la luz con gran precisión}
medir la velocidad en distintos puntos con gran precisión (osea la primera idea que teníamos) usando los métodos que en el pasado han dado resultados precisios: ramdani ref 3
\section{Circuito LC en el espacio}
Llevar el equipo descrito en el capitulo III a la estación espacial china.
\section{ Circuito LC con fem aplicada por inducción}

% \begin{figure}[h!]\centering\includegraphics[width=0.6\linewidth]{montaje2.png}\caption{Esquema del segundo montaje.}\end{figure}

En esta segunda forma se conectan la bobina y el condensador en serie, luego se enrolla un cable no aislado alrededor del inductor y los finales de este a un generador de frecuencia, de igual modo se usa otro cable para enrollarlo alrededor del condensador y se conectan sus extremos a un osciloscopio. Este montaje está basado en la referencia \cite{GuiaLAB-MIT}. En este caso la corriente variable $i_0(t)$ induce una fem $\varepsilon$ en el circuito (ley de Faraday) que varía con la misma frecuencia que la corriente original, a su vez $\varepsilon$ genera una corriente inducida $i(t)$ en el circuito LC que tiene la misma frecuencia que $i_0(t)$. Para realizar las mediciones, en la misma referencia [2] se plantea enrollar otro cable no aislado alrededor del inductor (en otra zona del mismo como se ve en la siguiente imagen) y conectar los extremos de este a un osciloscopio.